\section{Complete complex powers on local time charts}
\label{sec:complex}

The rational moments are regular at a merger. To extend the denominator
power to a complex index, the logarithm must be chosen on a spatial
neighborhood of each coalescing cluster. The two infinite ends of the
divisor admit a common principal branch. We use the stated domain
\eqref{eq:domain}, the strip and estimates \eqref{eq:count-safe}, and the
frozen \(\ell\) throughout this section \cite{BakerBase,BakerGT41,BakerCM03}.

Use \(w_j=\rho_j-c\), \(Q_j=Q(t,w_j)\) and
\(\mu_j=\ell^2/Q_j\) from \eqref{eq:quadratic}, with every exterior
occurrence counted. The collar gives \(|Q_j|\ge(15/16)|w_j|^2\).
For a specified spatial logarithm of \(Q(t,w)\), put
\[
 f_p(t,w)=\exp\!\big(p(2\log\ell-\Log Q(t,w))\big),\qquad
 h(t,w)=\frac{H_z(t,c(t)+w)}{H(t,c(t)+w)} .
\]
Here \(\log\ell\) is real. In a time chart, define
\(Z(p,t)=\sum_j f_p(t,w_j)\). At positive integers \(n\) this agrees
with \(\ell^{2n} S_n(t)\).

\subsection{Spatial contours and the time law}

At each reference time choose left and right centered cuts
\(-R_-\in[-7,-6]\) and \(R_+\in[6,7]\) with no zero on their
vertical connectors \(|\Im w|\le4\). These cuts remain cleared on
a smaller real-time interval. Appendix~\ref{app:clearance} proves the
choice from the unit count and gives the bounds needed to remove remote
vertical closures. On the two closed half-strips beyond the cuts,
\(\Re Q\ge(\Re w)^2/2\), so we use the principal log. This defines
the power at every exterior occurrence beyond either cut.

Inside the central rectangle, enclose each distinct reference
exterior location in a small disk. The disks have disjoint closures
and avoid the target and \(Q=0\). Their boundaries \(\gamma_C\) are
zero-free. The argument principle preserves each count and the total
central count on a smaller interval, so the disks account for all
central exterior occurrences. At a disk center \(w_*\) and reference
time \(t_*\) choose \(L_*\) with \(\exp L_*=Q_*=Q(t_*,w_*)\) and
\(\Im L_*\in[-\pi,\pi]\). On the disk use
\begin{equation}
 \Log Q=L_*+\Log_{\rm pr}(Q/Q_*),\qquad |Q/Q_*-1|<1/2 .
 \label{cp:local-log}
\end{equation}
Shrinking the disks and the time interval ensures the inequality on
neighborhoods of their closures. These logs have
\(|\Im\Log Q|\le\Theta=3\pi/2\). Each disk uses one log through any
merger of its counted occurrences. The compact established interval admits a
finite cover by such charts, with estimates on compact subintervals
that still cover it.

For an integrand \(B\) with the log prescribed on each piece, let
\begin{equation}
\begin{aligned}
 \mathcal S(B)=\frac1{2\pi\ii}\Bigg(&
   \sum_C\oint_{\gamma_C} B(w)\dd w
   +\int_{R_+}^{\infty}\!\big[B(x-4\ii)-B(x+4\ii)\big]\dd x\\
 &+\int_{-\infty}^{-R_-}\!\big[B(x-4\ii)-B(x+4\ii)\big]\dd x\\
 &+\int_{R_++4\ii}^{R_+-4\ii} B(w)\dd w
  +\int_{-R_--4\ii}^{-R_-+4\ii} B(w)\dd w\Bigg).
\end{aligned}
\label{cp:spatial-boundary}
\end{equation}
The circles are counterclockwise. Both lower rays go rightward and
both upper rays leftward. The finite right connector goes down, and
the finite left connector goes up.

\begin{theorem}[Complete local complex-power law]\label{thm:complex-time}
On each chart the complete sum \(Z\) is holomorphic in \(p\) for
\(\Re p>1/2\) and is \(C^1\) in real physical time, with inward
derivatives at the interval endpoints. With derivatives at fixed \(w\),
\begin{equation}
\begin{aligned}
 Z&=\mathcal S(f_ph),\\
 Z_t&=\mathcal S((f_p)_t h+f_ph_t)
   =\mathcal S\big(\big[(f_p)_t-(f_p)_{ww}-c'(f_p)_w\big]h
                                       +(f_p)_wh^2\big),\\
 h_t&=-h_{ww}-2h h_w+c' h_w,\qquad (f_p)_t=pD'f_p/Q .
\end{aligned}
\label{eq:complex-time}
\end{equation}
The derivative is holomorphic on the same half-plane and continuous
through target and exterior mergers on the chart.
\end{theorem}

\begin{proof}
The complete value bound in Appendix~\ref{app:value-sums} gives normal
convergence in \(p\). For the spatial formula, first close the two
half-strips at distant cleared verticals, upward on the right and
downward on the left. The residue of \(h\) is the zero multiplicity,
and \(f_p\) is holomorphic in the region of each contour. Hence the
finite contours count all their exterior occurrences. The remote
edges vanish by \eqref{app:edge-bound}. The omitted divisor values
tend to zero and the horizontal rays converge absolutely.

The centered heat equation yields
\(h_t=\partial_w(-H_{zz}/H)+c'h_w\). Substituting
\(H_{zz}/H=h_w+h^2\) gives the last line of
\eqref{eq:complex-time}. Differentiation under \(\mathcal S\) is
justified independently of divisor truncation by the safe-line and
finite-contour estimates in Appendix~\ref{app:time-bounds}. For the
second derivative formula, the difference of the two integrands is
the spatial derivative of
\[
 -f_ph_w+(f_p)_wh-f_ph^2+c'f_ph .
\]
The integral of this derivative vanishes on each circle. On each
tail boundary the primitive agrees where the rays meet the finite
connector and tends to zero at the infinite ends. The square of
the full field \(h\) in \eqref{eq:complex-time} retains the
spatial interaction in the integer law \eqref{eq:low-orders}.
\end{proof}

Two time charts use the same principal logs for \(|\Re w|\ge7\).
Near each overlap time, refine the finitely many other assignments
by small disks around the distinct affected locations, with fixed
counts. In each refinement disk the two logs differ by a constant
\(2\pi\ii k_C\), \(k_C\in\mathbb{Z}\). The branch difference is
\begin{equation}
 Z_B-Z_A=\sum_C\big(\exp(-2\pi\ii k_Cp)-1\big)Z_{A,C},
 \label{cp:transition}
\end{equation}
where \(Z_{A,C}\) is the trace of the power in chart \(A\) over
that disk. Its fixed circle proves validity also after one
physical time differentiation. All these finite corrections
vanish at integers, including their time derivatives. This is a
local atlas for the noninteger powers.

\subsection{Conjugation and the negative-axis defect}

For a function of \(p\), write
\(Z^\sharp(p,t)=\overline{Z(\bar p,t)}\). The field and selected
pair are conjugation invariant at real time, so the central disks
can be chosen in conjugate pairs, with self-conjugate disks at
real centers. For real exterior \(w\) the denominator \(Q\) is
positive. A negative-real \(Q\) can arise only at \(w=\ii y\ne0\)
since \(\Im Q=2\Re w\,\Im w\). Its two conjugate locations are
distinct. Choose opposite arguments \(+\pi,-\pi\) there and take
conjugate \(L_*\) on all paired disks, with the real log at real
centers. In \eqref{cp:local-log} the ratio log has argument of
magnitude less than \(\pi/4\), so \(5\pi/4\) suffices for \(\Theta\)
on these charts. Both far principal logs are conjugation
compatible. Conjugation now permutes the entire occurrence
list and conjugates the logs. Absolute convergence gives
\begin{equation}
 Z^\sharp=Z
 \label{cp:schwarz}
\end{equation}
locally in real time, including mergers. This symmetry also holds
for \(Z_t\) by Theorem~\ref{thm:complex-time}.

At a fixed slice let \(Z_+\) have principal logs off the negative
axis and argument \(+\pi\) at all negative-real \(Q_j\). Its exact
defect is
\begin{equation}
 Z_+-Z_+^\sharp=-2\ii\,\ell^{2p}\sin(\pi p)
                        \sum_{j:Q_j<0}|Q_j|^{-p}.
 \label{cp:axis-defect}
\end{equation}
The sum is finite and includes both conjugate locations with
all multiplicities. This follows by subtracting the two
exponentials \(\exp(-\ii\pi p)\) and \(\exp(+\ii\pi p)\) at each
such occurrence. Resetting to this slice prescription can cross
a branch cut. The time law uses the continuations of the local logs.
At that fixed slice the algebra of \eqref{cp:transition} compares this
prescription with any chart: group the finitely many affected occurrences
by their logarithm shifts. This value comparison does not assert that
resetting to $Z_+$ defines an analytic time chart. The differentiated
transition formula applies to pairs of local analytic charts
\cite{BakerCM04,BakerCMReview}.

